\begin{textsample}{POS Dim 1 – df – Score 41.00 – df/df\_em\_15.txt}  \label{ex:f1_pos_129}
5.1 LINGUAGENS E SUAS \textbf{TECNOLOGIAS} - FORMAÇÃO GERAL BÁSICA As linguagens são sistemas de comunicação que tornam possíveis a organização e a manifestação do pensamento, o que favorece a \textbf{sistematização}, a difusão e o compartilhamento de conhecimentos, assim como a reflexão, a (re)criação e a (re)produção deles.
É por meio das linguagens que as novas gerações têm acesso ao conhecimento acumulado da \textbf{humanidade} e conseguem utilizá -lo como base para a construção de novos conceitos.
Além disso, as linguagens favorecem a fruição de saberes científicos, históricos, culturais, artísticos, estéticos, poéticos e de expressão humana, fortalecendo, assim, o modo de ser no mundo e de pertencer a ele.
Vale ressaltar que a nomenclatura “ Linguagens e suas Tecnologias ” está \textbf{baseada} nas atualizações da multiplicidade de linguagens, semioses e mídias envolvidas na criação de significação para os textos multimodais \textbf{contemporâneos} e, por outro \textbf{lado}, para a \textbf{pluralidade} e a diversidade cultural trazidas pelos autores/leitores \textbf{contemporâneos} a essa criação de significação ( ROJO, 2013 ).
Ademais, os \textbf{multiletramentos} estão relacionados com a \textbf{atual} produção dos textos, em seus diversos formatos e modalidades, e os meios em que eles estão alojados.
Sendo assim, os textos mudam seus formatos a partir da multimodalidade - \textbf{infográficos}, fluxogramas, gráficos, etc. ( RIBEIRO, 2016 ) –, e os leitores e produtores de textos também mudam em função da cultura global, que nesse caso é trazida pelas Tecnologias \textbf{Digitais} da Informação e Comunicação ( TDIC ).
As Linguagens e suas Tecnologias, vistas como área do conhecimento, estão previstas nas Diretrizes Curriculares para o Ensino Médio, que orientam a organização dos currículos escolares a partir da premissa de que “ as linguagens são \textbf{indispensáveis} para a constituição de conhecimentos e \textbf{competências} ” ( BRASIL, 1998 ).
Esse entendimento acerca das linguagens foi recepcionado pela Resolução nº 3, de 21 de novembro de 2018, que atualiza as DCNEM e ratifica que Linguagens e suas Tecnologias \textbf{configuram} uma das áreas do conhecimento que compõem a Formação Geral Básica ( BRASIL, 2018a ).
Na BNCC para o Ensino Médio, a área é apresentada como oportunidade de \textbf{aprofundamento} e consolidação das aprendizagens elencadas na BNCC para o Ensino Fundamental referentes aos componentes curriculares Língua Portuguesa, Arte, Educação Física e Língua \textbf{Inglesa}.
Além disso, a Base traz a perspectiva de que a formação escolar em Linguagens deve ser “ voltada a possibilitar uma participação mais plena dos jovens nas diferentes práticas socioculturais que envolvem o uso das linguagens ”, o que significa valorizar os conhecimentos linguísticos prévios dos jovens e as práticas de linguagem que fazem parte de seu cotidiano ( BRASIL, 2018a ).
Nesse sentido, o presente Currículo propõe que os conteúdos na área de Linguagens sejam trabalhados de modo que os estudantes possam “ vivenciar experiências significativas com práticas de linguagem em diferentes mídias ( impressa, \textbf{digital}, analógica ), situadas em campos de atuação social diversos ” ( \textbf{grifos} no original ), conforme preceituados na BNCC para o Ensino Médio ( campo da vida social, campo das práticas de estudo e pesquisa, campo jornalístico-midiático, campo de atuação na vida pública e campo artístico ) ( BRASIL, 2018a, p. 485 ).
Para \textbf{tanto}, é fundamental que o ambiente escolar incentive a interação dos jovens com situações-problema nas quais eles precisem tomar decisões lúcidas, \textbf{críticas}, conscientes e \textbf{reflexivas}, exercitando a defesa de \textbf{posicionamentos} que contribuam para a formação cidadã e para o desempenho de um papel social dinâmico e colaborativo.
Diante desse desafio, apresenta -se uma área de Linguagens com propensão interdisciplinar, centrada no desenvolvimento cognitivo e \textbf{socioemocional} dos estudantes e pautada pela formação de cidadãos conhecedores de seus deveres e direitos, preparados para a busca do pleno exercício de sua cidadania, com atitudes éticas que possibilitem uma leitura \textbf{crítica} e \textbf{transformadora} da sociedade.
Assim, a área foi organizada para oferecer aos estudantes práticas pedagógicas diversificadas que lhes permitam interagir com discursos \textbf{multissemióticos} ( manifestações visuais, sonoras, verbais e corporais ), a fim de que, por meio do exercício da reflexão, possam desenvolver aprendizagens significativas o bastante para conduzi -los a uma atuação social consciente, ética, produtiva e sustentável.
Outrossim, a área foi estruturada de modo que o trabalho pedagógico contemple as mais variadas práticas de linguagens, abrangendo, ainda, gêneros multimidiáticos, sobretudo aqueles relacionados à produção e difusão de conteúdos \textbf{digitais}, uma vez que a interação social, por meio de funcionalidades disponibilizadas pela internet, é largamente presente no cotidiano da sociedade \textbf{contemporânea}.
Esse ambiente \textbf{digital} possui cultura própria e suas práticas linguísticas e comunicativas são elementos \textbf{relevantes} na formação dos jovens, uma vez que \textbf{impactam} suas vidas.
Todavia, \textbf{devido} a um cenário não igualitário de acesso à internet e às TDIC, deve -se \textbf{salientar} que o trabalho pedagógico também pode ser oportunizado por \textbf{intermédio} de recursos analógicos, de \textbf{maneira} a contemplar os estudantes que não possuem acesso aos meios \textbf{digitais}.
Nessa perspectiva, os \textbf{multiletramentos} e o uso de metodologias ativas fazem parte do planejamento pedagógico da área de Linguagens e suas Tecnologias, já que o Currículo em Movimento do Distrito Federal valoriza os \textbf{letramentos} locais como instrumento de reafirmação dos laços de identidade com a comunidade e consigo mesmo.
Isso se faz sem \textbf{restringir} o acesso dos estudantes àqueles já consagrados na sociedade ( especialmente a local ), a fim de oferecer -lhes a oportunidade de ampliar sua visão de mundo e suas possibilidades de interação sociocultural, proporcionando uma formação holística do indivíduo.
Diante do \textbf{exposto}, o Currículo em Movimento do Distrito Federal apresenta cinco objetivos gerais para a área de Linguagens : - favorecer a \textbf{compreensão} de diversificadas práticas de linguagem ( \textbf{multissemióticas} e multimidiáticas ) e práticas culturais ( verbais, artísticas e corporais ), promovendo a \textbf{mobilização} desses conhecimentos para o exercício \textbf{reflexivo} acerca de conteúdos informativos amplamente \textbf{divulgados} em mídias diversas, a fim de contribuir com a formação de um cidadão \textbf{crítico}, consciente, ético e protagonista, com elevado senso de identidade e de pertencimento social ; - promover o entendimento da importância da apropriação das práticas de linguagem, por meio da interação \textbf{crítica} com variadas manifestações linguísticas, corporais, artísticas e culturais, a fim de compreender a realidade e ampliar as possibilidades de atuação social, pautando -se pelos Direitos Humanos e pelos \textbf{ideais} de justiça social e democracia, repudiando quaisquer formas de preconceito ; - valorizar as diversas linguagens, abordando formas locais de expressão, estimulando o senso de pertença, a iniciativa e a \textbf{autoria} ( respeitando -se a originalidade do texto ), como instrumentos de apropriação da produção cultural coletiva e de participação social ; e, ainda, promovendo a \textbf{compreensão} de cânones regionais, nacionais e mundiais de expressão, a fim de expandir as possibilidades de interação cultural ; - propiciar vivências de práticas corporais, estéticas e linguísticas que possibilitem a \textbf{compreensão} das variedades culturais em níveis local, regional, nacional e mundial, de modo a construir múltiplas possibilidades de interação com o conhecimento acumulado da \textbf{humanidade}, com a finalidade de compreender a relevância das práticas de linguagens para as diversas culturas e construir relações privadas de preconceitos ; - apreciar as manifestações linguísticas, estéticas e corporais presentes na cultura \textbf{digital}, analisando as potencialidades do meio \textbf{digital} para a produção e a divulgação de informações, exercitando a \textbf{crítica} aos conteúdos disponibilizados pelas mídias, a fim de interagir, de forma produtiva, com os conhecimentos publicizados digitalmente e estimular a contribuição autoral, dentro dos princípios da legalidade de \textbf{autoria}.
Esses objetivos gerais dão origem a objetivos específicos que visam ampliar as possibilidades de promover uma formação integral e centrada no desenvolvimento de \textbf{competências} cognitivas e \textbf{socioemocionais} que favoreçam uma atuação autônoma, protagonista e ética dos estudantes em seu cotidiano.
Para \textbf{tanto}, os objetivos específicos foram agrupados em seis unidades temáticas, elencadas a seguir, cujo propósito é oferecer diretrizes quanto ao enfoque de formação humana a ser priorizado no planejamento e no fazer didático em cada semestre letivo : - Linguagens e suas \textbf{tecnologias} em contextos e práticas culturais ; - Linguagens e suas \textbf{tecnologias} em contextos e práticas sociais ; - Linguagens e suas \textbf{tecnologias} em e para contextos de Direitos Humanos ; - Linguagens e suas \textbf{tecnologias} em contextos \textbf{socioambientais} ; - Linguagens e suas \textbf{tecnologias} em contextos de identidade e protagonismo juvenil ; - Linguagens e suas \textbf{tecnologias} em contextos de cultura \textbf{digital}.
Além dessas seis unidades temáticas, este Currículo traz a perspectiva de divisão dos objetivos de aprendizagem específicos de Língua Portuguesa, \textbf{aproximando} -o da organização presente na BNCC.
Tal divisão se deu para que os objetos de conhecimento referentes ao ensino da língua materna fossem abordados pelos professores especialistas deste componente curricular, embora na área de Linguagens e suas Tecnologias existam objetivos de aprendizagem contendo objetos de conhecimento de Língua Portuguesa interdisciplinares com os demais componentes curriculares da área.
Para contemplar a formação integral dos estudantes do Ensino Médio, o Currículo em Movimento vislumbra uma atuação interdisciplinar, integrada, coordenada e colaborativa dos seguintes componentes curriculares :

% matched lemmas: aprofundamento, aproximar, atual, autoria, basear, competência, compreensão, configurar, contemporâneo, crítico, devido, digital, divulgar, exposto, grifo, humanidade, ideal, impactar, indispensável, infográfico, inglês, intermédio, lado, letramento, maneira, mobilização, multiletramento, multissemiótico, pluralidade, posicionamento, reflexivo, relevante, restringir, salientar, sistematização, socioambiental, socioemocional, tanto, tecnologia, transformador
\end{textsample}
